\chapter{Simulation}

The \name{} library contains some basic modules allowing simple discrete events simulations experiments.
This module has not the ambition to compete with commercial or well-established open-source packages, but to provide support tools for teaching and research.

\section{Introduction to discrete events simulation}

We consider a system characterized by the state $\mathcal{S}(t)$, where $t$ denotes the time.
In the discrete events setting, the system state can only change at a countable number of points in time.
The state can be modified when (and only when) some event happens, which is assumed to be instantaneous.
The system evolution is therefore the consequence of a sequence of
events, $e_0, e_1, e_2, \dots$, occuring at times $0=t_0 \le t_1 \le
t_2 \le \cdots$.
Let ${\mathcal{S}_i}$ the system state directly after ${e_i}$, and define
the simulation time as the current value of ${t_i}$.
We can represente the discrete events simulation as illustrated on
Figure~\ref{fig:discrete}.

\begin{figure}[htb]
\begin{center}
%%  Evolution of a general discrete-event model.
\hspace{0.5cm}
\beginpicture
\setcoordinatesystem units <7pt,7pt>
%%  72.27pt = 1in
\setplotarea x from 0 to 36, y from 0 to 10
\axis left
  label {\lines {\emph{state}\cr\vbox to 95pt{\hbox to 1.5pt{\null }} \cr}
         \hskip 9pt } /

\axis bottom
  label {\hbox to 4.0in {\hfill \emph{time}}}
  ticks length <4pt> withvalues 
     0 $t_1$ $t_2$ $t_3$ $t_4$ $t_5$ $t_6$ /
  at 0.0 6.0 7.5 12.5 18.0 25.2 30.0 / /

\axis bottom /
\put {$\rightarrow$} at 36.0 -0.1 
\put {$\uparrow$}    at 0.0 10.0 
\multiput {$\bullet$} at
  0.0 2.0
  6.0 1.6
  7.5 2.7
 12.5 6.1
 18.0 5.1
 25.2 2.6
 30.0 4.0 /
\putrule from  0.0 2.0 to  6.0 2.0
\putrule from  6.0 1.6 to  7.5 1.6
\putrule from  7.5 2.7 to 12.5 2.7
\putrule from 12.5 6.1 to 18.0 6.1
\putrule from 18.0 5.1 to 25.2 5.1
\putrule from 25.2 2.6 to 30.0 2.6
\putrule from 30.0 4.0 to 32.0 4.0 
\endpicture
\end{center}
\caption{}
\label{fig:discrete}
\end{figure}

The core of a discrete events simulation as proposed in \name{} is the \verb|ESim| structure,
which will be responsible of the simulation clock and events list gestion.
A new instance can be created with the function
\begin{verbatim}
ESim *esim_new()
\end{verbatim}
and suppressed with the function
\begin{verbatim}
void esim_free(Esim *sim);
\end{verbatim}

The second basic component of a discrete events simulation is an
event, that can be created with the function
\begin{verbatim}
ESim_Event *esim_event_new(ESim_Action action);
\end{verbatim}
and can be freed with the function
\begin{verbatim}
void esim_event_free(void *ev);
\end{verbatim}
The use of the \verb|void *| type makes the function compatible with
the utilities associated with the gestion of lists.

Since an event usually modifies the system state, it has to perform
some action.
The type \verb|ESim_Action| is a pointer to a function of the form
\begin{verbatim}
int (* ESim_Action)(void *system);
\end{verbatim}
The choice to use C instead of an object-oriented programming approach imply some technical limitations if we want to maintain generality.
The event function prototype considers a system as a \verb|void *| pointer, so that every implementable system can be considered.
It is the responsability of the programmer to correctly cast the \verb|void *| pointer in the events gestion functions.
